\subsection{HOMEOMORPHISMS OF OPEN INTERVALS OF R}

PROPOSITION 1. All non-empty open intervals of \(\mathbf{R}\) are homeomorphic to \(\mathbf{R}\) . Consider first a bounded open interval \(\mathbf{I} = ]a\) , \(b[(a < b)\) . For each \(x \in \mathbf{I}\) put \(f(x) = -\left(\frac{1}{x - a} + \frac{1}{x - b}\right)\) . This function is continuous and strictly increasing on \(\mathbf{I}\) , for we have seen that \(\frac{1}{x - b}\) is strictly decreasing in \(]\leftarrow, b[\) and \(\frac{1}{x - a}\) strictly decreasing in \(]a, \rightarrow[\) . It follows that \(f\) is a homeomorphism of \(\mathbf{I}\) onto an interval \(f(\mathbf{I})\) of \(\mathbf{R}\) ( \(\S 2\) , no. 6, Theorem 5). \(f(\mathbf{I})\) is neither bounded above nor below; if, for example, we had \(f(x) \leqslant c\) for all \(x \in \mathbf{I}\) , it would follow that \(1 - \frac{b - x}{x - a} \leqslant c(b - x)\) , since \(b - x > 0\) ; and this leads to a contradiction when \(x\) is sufficiently near \(b\) (by virtue of the continuity of the two sides of the inequality, which are rational functions, at the point \(b\) ). Hence \(f(\mathbf{I}) = \mathbf{R}\) , and therefore every bounded open interval is homeomorphic to \(\mathbf{R}\) . Let \(g\) be the inverse of \(f\) : it maps every unbounded open interval of \(\mathbf{R}\) onto an interval \(J\) contained in \(I\) and open in \(I\) . Since \(I\) is open in \(\mathbf{R}\) , \(J\) is also open in \(\mathbf{R}\) ; since it is bounded, it is homeomorphic to \(\mathbf{R}\) ; and thus we have proved that every unbounded open interval is homeomorphic to \(\mathbf{R}\) .

Remark. To show that all bounded open intervals are homeomorphic to each other, it would be enough to remark that, if \(a \neq b\) and \(a' \neq b'\) , there is a homeomorphism of R onto itself, of the form \(x \to \alpha x + \beta\) (and only one) which maps a to \(a'\) and b to \(b'\) , and hence maps the open (resp. half-open, closed) interval with end-points a, b onto the open (resp. half-open, closed) interval with end-points \(a', b'\) ; the reader may easily verify this by calculating \(\alpha\) and \(\beta\) .

